\loai{Xác định cảm ứng từ}
\begin{vidu} %[2P4-2.2B][Loại1]
Một dây dẫn điện thẳng dài có dòng điện chạy qua với cường độ $I=0,5$ A đặt trong không khí. Tính cảm ứng từ đặt tại M cách dòng điện $r=4\ cm$.
\choice
{\True $0,25.10^{-5}\ T$}
{$0,35.10^{-5}\ T$}
{$0,45.10^{-5}\ T$}
{$0,55.10^{-5}\ T$}	
\loigiai{Cảm ứng từ tại M: $B=2.10^{-7}\dfrac{I}{r}=2.10^{-7}\dfrac{0,5}{0,04}=0,25.10^{-5}\ T$.
}
\end{vidu}
\loai{Xác định cường độ dòng điện}
\begin{vidu} %[2P4-2.2B][Loại2]
Từ cảm B của dòng điện thẳng tại điểm M cách dòng điện 3 cm bằng $2,4.10^{-5}$ T. Tính cường độ dòng điện của dây dẫn.
\choice
{0,36 A}
{0,72 A}
{\True 3,6 A}
{7,2 A}
\loigiai{Dòng điện thẳng: $B=2.10^{-7}.\dfrac{I}{r}\Rightarrow I=\dfrac{B\text{r}}{{2.10}^{-7}}=3,6\ A$}
\end{vidu}
\loai{Xác định khoảng cách}
\begin{vidu} %[2P4-2.2B][Loại3] 
Dây dẫn thẳng dài có dòng điện 5 A chạy qua. Cảm ứng từ tại M có độ lớn $10^{-5}\ T$. Điểm M cách dây một khoảng
\choice
{20 cm}
{\True 10 cm}
{5 cm}
{2 cm}
\loigiai{Cảm ứng từ gây bởi dòng điện thẳng dài được xác định bằng biểu thức:\\ $B=2.10^{-7}\dfrac{I}{r}\Rightarrow r=\dfrac{2.10^{-7}I}{B}=\dfrac{2.10^{-7}.5}{10^{-5}}=10\ cm$}
\end{vidu}
\loai{Mối liên hệ giữa cảm ứng từ với khoảng cách}
\begin{vidu} %[2P4-2.2B][Loại4]
Một điểm cách một dây dẫn dài vô hạn mang dòng điện 20 cm thì có độ lớn cảm ứng từ $1,2\ \mu T$. Một điểm cách dây dẫn đó 60 cm thì có độ lớn cảm ứng từ là
\choice
{\True $0,4\ \mu T$}
{$0,2\ \mu T$}
{$3,6\ \mu T$}
{$4,8\ \mu T$}
\loigiai{Ta có: $B=2.10^{-7}.\dfrac{I}{r}$ nên $\dfrac{B_1}{B_2}=\dfrac{r_2}{r_1}=3\Rightarrow B_2=\dfrac{B_1}{3}=0,4\ \mu T$}
\end{vidu}
\begin{vidu} %[2P4-2.2K][Loại4]
Tại một điểm cách một dây dẫn thẳng dài vô hạn mang dòng điện 5 A cảm ứng từ 0,4 $\mu T$. Nếu cường độ dòng điện trong dây dẫn tăng thêm 10 A cảm ứng từ tại điểm đó có giá trị là
\choice
{1,6 $\mu T$}
{0,2 $\mu T$}
{\True 1,2 $\mu T$}
{0,8 $\mu T$}
\loigiai{Ta có $B = 2.10^{- 7}\dfrac{I}{r} \Rightarrow$ khi $I_2 = 3I_1 \Rightarrow B_2= 3B_1$}
\end{vidu}

